
@@M0@@ le nombre d'intersection, qui peut être calculé comme coefficient du monôme @@M1@@ dans le polynôme de Snapper @@M2@@ (X 4.3.1).

\textbf{2.2.} Désormais, on suppose que @@M3@@ est algébriquement clos, que @@M4@@ est projectif, muni d'un faisceau inversible très ample @@M5@@ et que @@M6@@ est intègre, de dimension @@M7@@.

Soit @@M8@@ un faisceau inversible sur @@M9@@. L'entier @@M10@@ est appelé le \uline{degré} de @@M11@@ (relativement à @@M12@@) et est noté @@M13@@.

Dans la suite on considère un faisceau inversible @@M14@@ et les polynômes de Hilbert : @@M15@@ et @@M16@@.

\textbf{Lemme 2.3.} \emph{On a} :
\begin{enumerate}[label=(\roman*)]
\item \emph{@@M17@@.}
\item \emph{(Si @@M18@@), @@M19@@.}
\end{enumerate}

En effet, on obtient (i) et (ii) aussitôt en faisant @@M20@@ dans le polynôme suivant :
@@M21@@

\textbf{Lemme 2.4.} \emph{Si @@M22@@ et @@M23@@, alors @@M24@@; donc @@M25@@ et @@M26@@.}

En effet, comme @@M27@@ est intègre, il existe un diviseur @@M28@@ tel que @@M29@@. Si @@M30@@, on aurait @@M31@@.

\textbf{Lemme 2.5.} \emph{Si @@M32@@, alors @@M33@@ est un @@M34@@-faisceau (avec @@M35@@ zéros).}

